\section{The endpoint mechanism}\label{sec:endpoint}

Let $k$ be the original coefficient field, and suppose
\[
 q\ge2,\qquad 0<a_1<\cdots<a_q<u<m,\qquad d\ge1,
 \qquad \gcd(a_1,\ldots,a_q,u)=\gcd(m,d)=1.
\]
Set
\[
 \Gamma=\langle m,m+da_1,\ldots,m+da_q,m+du\rangle,
 \qquad G=\operatorname{gr}_{\mathfrak m}k[[t^\Gamma]],
\]
and write $S=k[x,y_1,\ldots,y_q,v]$, $E=k[x,v]$ and
$T=k[y_1,\ldots,y_q]$, with every variable of degree one.
We use $\beta_i(G)$ over $S$ and $b_i(A)$ over $T$.
Binomial coefficients and Betti numbers with negative subscripts are zero.
The physical width is $V=du$; $u$ is the normalized width.
In particular, normalization retains $m$ and $d$ and describes the same
semigroup and tangent cone. The endpoint construction needs $u<m$;
the narrow-width application also has $V\le m-2$.

\paragraph{The core and its fibers.}
Consider the standard graded algebra
\[
 C=k[s,st^{a_1},\ldots,st^{a_q},st^u],
 \qquad x=s,\quad y_j=st^{a_j},\quad v=st^u.
\]
The displayed generators of $\Gamma$ are minimal. Indeed, a decomposition
of the generator with offset $a\in\{0,a_1,\ldots,a_q,u\}$ into
$l\ge2$ generators, with offset sum $B$, would give
$(l-1)m=d(a-B)$. Coprimality forces $l-1=jd$ and $a-B=jm$
for some $j\ge1$, contradicting $a\le u<m$.
If $\nu(h)$ denotes the maximum length of a generator expression for $h$,
there is a graded surjection
\begin{equation}\label{eq:endpoint-core-map}
 C\longrightarrow G,\qquad
 s^nt^B\longmapsto
 \begin{cases}(t^{nm+dB})^*,&\nu(nm+dB)=n,\\
 0,&\nu(nm+dB)>n.
 \end{cases}
\end{equation}
Maximum-length expressions prove surjectivity and compatibility with products.
Distinct offsets in a fixed degree give distinct numerical exponents, so
its kernel is spanned by the core monomials that die.

For $0\le r<u$, put $\epsilon_r=0$ if $r=0$ and $\epsilon_r=1$ otherwise.
Write the degree and offset of a core monomial as
\[
 B=r+uK,\qquad n=L+K+\epsilon_r.
\]
Since $B\le nu$, both $L$ and $K$ are nonnegative.
Multiplication by $x$ or $v$ increases respectively $L$ or $K$.
Let $M_r\subseteq\N^2$ be the resulting upward closed set, also viewed
as a monomial ideal of $E$. Thus, as $E$-modules,
\begin{equation}\label{eq:endpoint-fibers}
 C\simeq\bigoplus_{r=0}^{u-1}M_r,
\end{equation}
with the degree of $(L,K)$ in the $r$th summand equal to
$L+K+\epsilon_r$.

These ideals have a finite, intrinsic description. A word of at least $u$
internal letters has two prefix sums equal modulo $u$. The intervening
nonempty word has length $l$ and offset $cu$, where $1\le c<l$;
replace it by $x^{l-c}v^c$. Repetition leaves at most $u-1$ internal
letters. Consequently every minimal point of $M_r$ is represented by a
pure internal word of length at most $u-1$. All residues occur, because
the internal offsets generate $\mathbb Z/u\mathbb Z$ and inverses in this
finite group can be expressed by nonnegative repetitions. Moreover,
$M_0=E$, and every other minimal point satisfies $L+K\le u-2$.
In particular $C$ is finite over $E$, and
\[
 H_x:=\max\{L:(L,K)\in\operatorname{Min}M_r\text{ for some }r\}
 \le u-2.
\]

\paragraph{All competitors and finite generation.}
Equality of numerical exponents shows that a longer expression for a
core monomial $(n,B)$ is exactly a core point
\[
 (n+jd,B-jm),\qquad j\ge1.
\]
Fix the longer-side residue $r$ and a minimal point $(L,K)$ of $M_r$.
Write $r+jm=s+cu$, with $0\le s<u$. On the shorter side this
competitor forbids the quadrant starting at
\begin{equation}\label{eq:endpoint-competitor}
 (A,B')=
 \bigl(\max\{0,L-jd-c+\epsilon_r-\epsilon_s\},K+c\bigr).
\end{equation}
The maximum accounts exactly for increasing the longer-side $x$ exponent
until the shorter-side first coordinate is nonnegative.
Let $T_s\subset E$ be generated by these starts for every $j\ge1$,
every $r$, and every minimal point of $M_r$. Then
\begin{equation}\label{eq:endpoint-kernel}
 N_s=M_s\cap T_s,\qquad
 G\simeq\bigoplus_{s=0}^{u-1}M_s/(M_s\cap T_s)
\end{equation}
as $E$-modules. The intersection imposes membership in the core;
the competitor ideal alone need not be the kernel.

Although all positive $j$ enter this description, the generating process
is finite. Indeed,
\[
 c-\epsilon_r+\epsilon_s
 =\left\lceil\frac{r+jm}{u}\right\rceil
  -\left\lceil\frac r u\right\rceil
 \ge\left\lfloor\frac{jm}{u}\right\rfloor\ge j,
 \qquad
 A\le\max\{H_x-2j,0\}.
\]
Thus all starts have first coordinate at most $H=\max\{H_x-2,0\}$.
For each first coordinate that occurs, retain only the smallest second
coordinate; these finitely many starts generate $T_s$.
Choosing $j$ with $2j\ge H_x$ and $r\equiv s-jm\pmod u$ also gives a
start with first coordinate zero. Hence every $T_s$ contains a pure
$v$ power. This proves the finite description directly from all
competitors, without a coprimality assumption on $m$ and $u$.

\paragraph{Length and socle resources.}
Define
\[
 \begin{gathered}
 Q_C=C/(x,v)C,\qquad Q_G=G/(x,v)G,\\
 P_C=\dim_k Q_C,\quad P_G=\dim_k Q_G,\quad
 R=\dim_k(0:_G(x,v)).
 \end{gathered}
\]
The minimal fiber points form a basis of $Q_C$, and their surviving
subset forms a basis of $Q_G$. As a $T$-algebra $Q_C$ is standard graded
Artinian, with Hilbert function
$h_n=\dim_k(Q_C)_n$ satisfying $h_0=1$, $h_1=q$ and $h_n=0$ for
$n\ge u$. It need not be a monomial quotient of $T$.
Let $\ell_r$ be the shortest internal word length reaching $r$ modulo $u$,
with $\ell_0=0$, and put
\begin{equation}\label{eq:endpoint-resources}
 W=1+\sum_r\ell_r,\qquad
 \widehat R=\min\{W-u,u\max\{H_x-2,0\}\}.
 \qquad P_C\le W,\quad R\le\widehat R.
\end{equation}

Here is the counting argument. A shortest word determines a minimal
point $(L_r,K_r)$. In the ordered staircase of minimal points, at most
$L_r$ points lie to its left and at most $K_r$ to its right. For
$r\ne0$ this gives at most $L_r+K_r+1=\ell_r$ points; residue zero
contributes one. This proves $P_C\le W$ and also $W\ge u$.

Order the minimal generators of $T_s$ as
\[
 (A_0,B_0),\ldots,(A_z,B_z),\qquad
 0=A_0<\cdots<A_z,\quad B_0>\cdots>B_z=\kappa_s.
\]
The socle of $E/T_s$ has the $z=z_s$ corner monomials
$(A_{j+1}-1,B_j-1)$ as a basis. This assertion only concerns its socle;
$E/T_s$ need not have finite length. Counting generators on either side
of a minimal point, and then passing to any dominating point, gives
\[
 z_s\le A+B-\kappa_s\quad((A,B)\in T_s),\qquad
 \kappa_s\ge\left\lceil\frac{m-s}{u}\right\rceil.
\]
The natural injection $M_s/(M_s\cap T_s)\hookrightarrow E/T_s$ induces
a socle injection, so $R\le\sum_s z_s$.
Apply the $j=1$ competitor to each shortest point, and set
$s_r=(r+m)\bmod u$, $c_r=(r+m-s_r)/u$.
Its start has $A\le L_r$ and $B=K_r+c_r$.
Translation permutes the residues, and
\[
 \sum_r c_r=m=\sum_s\left\lceil\frac{m-s}{u}\right\rceil.
\]
Summing the point bounds therefore yields
$\sum_s z_s\le\sum_r(L_r+K_r)=W-u$.
Independently, the first coordinates of minimal generators lie in
$[0,H]$, whence $z_s\le H$. Together these prove
\eqref{eq:endpoint-resources}. Every resource refers to the same
orientation of the core; replacing $H_x$ by a reflected height would
require a separate comparison with this $R$.

\paragraph{The compensated Tor bound.}
The endpoint Koszul homology is finite-dimensional and has dimensions
\begin{equation}\label{eq:endpoint-homology}
 \bigl(\dim_k H_0(x,v;G),\dim_k H_1(x,v;G),
       \dim_k H_2(x,v;G)\bigr)=(P_G,P_G+R,R).
\end{equation}
For completeness, $G/xG$ and $(0:_Gx)$ have finite dimension by the
Ap\'ery description of the multiplicity action. View
$K(x,v;G)$ as the cone of multiplication by $v$ on $K(x;G)$.
The cone long exact sequence, applied to these two finite homology
groups, gives the middle equality in \eqref{eq:endpoint-homology} by
rank--nullity. The zeroth and second groups are respectively $Q_G$ and
$(0:_G(x,v))$. No finite-dimensionality assertion about $G$ is involved.

Let $K$ be the kernel of the actual $T$-linear surjection $Q_C\to Q_G$.
Then $\dim_k K=P_C-P_G\ge0$. The Tor long exact sequence and the
dimensions of the internal Koszul chain groups give
\[
 b_i(Q_G)\le b_i(Q_C)+(P_C-P_G)\binom q{i-1}.
\]
In the bounded double Koszul complex on the actual $G$, take endpoint
homology first. Its zeroth row retains the genuine internal Tor of
$Q_G$; bounding the other two rows by their chain dimensions gives
\[
 \beta_i(G)\le b_i(Q_G)+(P_G+R)\binom q{i-1}
                              +R\binom q{i-2}.
\]
This is a dimension inequality for a convergent filtration, and requires
no degeneration assertion. Substitution cancels $P_G$ against the
quotient-kernel length. Pascal's identity yields
\begin{equation}\label{eq:endpoint-bound}
 \boxed{\displaystyle
 \beta_i(G)\le b_i(Q_C)+P_C\binom q{i-1}
                              +R\binom{q+1}{i-1}.}
\end{equation}
In particular $\beta_0(G)=1$, $\beta_{q+2}(G)\le R$, and
$\beta_i(G)=0$ for $i>q+2$.

\paragraph{Artinian estimates and the exported rows.}
For a standard graded Artinian $A=T/I$ with $I_1=0$, write
$P=\dim_k A$ and $h_n=\dim_k A_n$.
The cyclic, degree-zero generated case of
Lemma~\ref{lem:ordinary-harmonic} gives
\begin{align}
 b_1(A)&\le q+\sum_{n\ge1}\frac{(q-1)h_n}{n+1}
 \le \frac{q-1}{3}P+\frac{(q+1)(q+2)}6,
 \label{eq:endpoint-artin-first}\\
 b_1(A)&\le\frac{q-1}{4}P+\frac{(q+1)^2}{4}
                         +\frac{q-1}{12}h_2.\nonumber
\end{align}
To see the specialization, the first partial Euler sum is
$c_1(A)=b_1(A)-b_0(A)=b_1(A)-1$.
The lemma bounds it by
$(q-1)\sum_{n\ge0}h_n/(n+1)$, and $h_0=1$ supplies the constant
$q$ in \eqref{eq:endpoint-artin-first}. Its hypotheses hold because
the class of $1$ generates $A$ in degree zero and the homogeneous
pieces have finite support.
The second and third bounds separate degree one, or degrees one and two,
and bound the remaining reciprocal weights by $1/3$ or $1/4$.
Subtracting $P$ from the harmonic bound gives
\[
 b_1(A)-P\le q-1+\sum_{n\ge1}\frac{q-n-2}{n+1}h_n.
\]
For $q=5$, discard the nonpositive terms with $n\ge3$ and use
$h_2\le15$ to obtain $b_1(A)-P\le14$.
All these weights and signed differences are interpreted in $\mathbb Q$.

There is also a uniform estimate for every positive row:
\[
 b_i(A)+\binom{q+1}{i+1}\le P\binom qi\qquad(i\ge1).
\]
Indeed, in the Koszul differential $d_i$, restrict the source to
coefficient degree zero. For each wedge $J$, select the target coordinate
$(y_{\min J},J\setminus\{\min J\})$. These coordinates give a diagonal
minor with entries $\pm1$, since the variable images in $A_1$ are
independent. Hence $\operatorname{rank}d_i\ge\binom qi$; apply the
same argument to $d_{i+1}$ and subtract both ranks from the chain
dimension $P\binom qi$.
The minors remain invertible in every characteristic, including two.

Applying these estimates to $Q_C$ in \eqref{eq:endpoint-bound} gives
\begin{equation}\label{eq:endpoint-rows}
 \boxed{\begin{aligned}
 \beta_1(G)&\le P_C+
  \left\lfloor q+\sum_{n\ge1}\frac{(q-1)h_n}{n+1}\right\rfloor
  +\widehat R,\\
 \beta_i(G)+\binom{q+1}{i+1}
 &\le P_C\binom{q+1}i+\widehat R\binom{q+1}{i-1}
 \qquad(i\ge1).
 \end{aligned}}
\end{equation}
The floor is valid because the Betti number is integral.
These estimates concern the original completed tangent cone over $k$;
no extension of the coefficient field is needed here.

Finally suppose $q=2$, with internal offsets $a<b$ and variables $y,z$.
In fixed degree $n$, the words $y^{n-j}z^j$ have distinct offsets
$na+j(b-a)$, so $Q_C$ is a monomial quotient of $k[y,z]$.
Put $\sigma_C=\dim_k(0:_{Q_C}(y,z))$.
The two-variable Koszul complex and its Euler identity give
$(b_0(Q_C),b_1(Q_C),b_2(Q_C))=(1,\sigma_C+1,\sigma_C)$.
Thus \eqref{eq:endpoint-bound} specializes to
\begin{equation}\label{eq:endpoint-two-rows}
 \boxed{\begin{aligned}
 \beta_1(G)&\le P_C+\sigma_C+1+R,&
 \beta_2(G)&\le2P_C+\sigma_C+3R,\\
 \beta_3(G)&\le P_C+3R,&\beta_4(G)&\le R.
 \end{aligned}}
\end{equation}
Each $R$ may be replaced by $\widehat R$. The internal socle
$\sigma_C$ and the endpoint socle $R$ belong to different quotients
and different variable actions.

The core, competitor, resource and transfer arguments also apply with
one internal offset. For $q=1$, coprimality says $\gcd(a_1,u)=1$.
The relation $y^u=x^{u-a_1}v^{a_1}$ in $C$ shows that $Q_C$ is
spanned by the classes of $1,y,\ldots,y^{u-1}$, so $P_C\le u$.
The one-variable Artinian presentation is
$Q_C\simeq k[y]/(y^{P_C})$, whose Betti numbers are $(1,1)$.
Formula \eqref{eq:endpoint-bound} therefore gives
\[
 (\beta_1(G),\beta_2(G),\beta_3(G))
 \le (P_C+R+1,P_C+2R,R)
\]
coordinatewise. These rows use the same actual endpoint socle $R$;
any sharper bound for that socle can be substituted before comparing
with the interval target.
